\section*{Capítulo \ref{ch:g11:fossil-fuel-geology} --- La geología del carbón, el petróleo y el gas}

\begin{solution}{exo:g11:fossil-fuel-geology:1}
El plancton muerto se deposita en un fondo marino pobre en oxígeno; la
materia orgánica se convierte en \omterm{def:g11:fossil-fuel-geology:source-rock}{querógeno} en un fango; el enterramiento bajo
sedimentos más jóvenes; la \omterm{def:g11:fossil-fuel-geology:oil-window}{maduración} en la \omterm{def:g11:fossil-fuel-geology:oil-window}{ventana del petróleo}; la
\omterm{def:g11:fossil-fuel-geology:migration}{migración}; el petróleo se acumula bajo una \omterm{def:g11:fossil-fuel-geology:reservoir}{roca sello}.
\end{solution}

\begin{solution}{exo:g11:fossil-fuel-geology:2}
(1) \omterm{def:g11:fossil-fuel-geology:source-rock}{Roca generadora}. (2) \omterm{def:g11:fossil-fuel-geology:reservoir}{Roca almacén}. (3) \omterm{def:g11:fossil-fuel-geology:reservoir}{Roca sello}. (4) \omterm{def:g11:fossil-fuel-geology:reservoir}{Roca sello}. (5) \omterm{def:g11:fossil-fuel-geology:reservoir}{Roca almacén}.
\end{solution}

\begin{solution}{exo:g11:fossil-fuel-geology:3}
Techo: $(60 - 10)/25 = \qty{2.0}{km}$; base: $(120 - 10)/25 = \qty{4.4}{km}$.
La ventana tiene $60/25 = \qty{2.4}{km}$ de espesor.
\end{solution}

\begin{solution}{exo:g11:fossil-fuel-geology:4}
(1) Turba, lignito, \omterm{def:g11:fossil-fuel-geology:coal}{carbón} bituminoso, antracita. (2) La antracita, con un 86
a \qty{97}{\percent} de carbono. (3) La antracita, más allá de unos
\qty{5}{km} y \qty{150}{\degreeCelsius}.
\end{solution}

\begin{solution}{exo:g11:fossil-fuel-geology:5}
(1) Los fluidos se ordenan según su densidad: el gas es el más ligero, luego
el petróleo y luego el agua. (2) El petróleo y el gas ascienden porque son
más ligeros que el agua; solo los detiene una roca que no puedan atravesar.
La fangolita (o la lutita) y la sal.
\end{solution}

\begin{solution}{exo:g11:fossil-fuel-geology:6}
(1) A: $(60 - 15)/20 = \qty{2.25}{km}$ y $(120 - 15)/20 = \qty{5.25}{km}$.
B: $45/40 \approx \qty{1.13}{km}$ y $105/40 \approx \qty{2.63}{km}$.
(2) A \qty{2}{km}: en A, $15 + 40 = \qty{55}{\degreeCelsius}$, todavía no
madura, sin petróleo; en B, $15 + 80 = \qty{95}{\degreeCelsius}$, en la
\omterm{def:g11:fossil-fuel-geology:oil-window}{ventana del petróleo}, produce petróleo. (3) En la cuenca A: allí el petróleo
solo se forma por debajo de \qty{2.25}{km}.
\end{solution}

\begin{solution}{exo:g11:fossil-fuel-geology:7}
(1) $11 \times 3 = \qty{33}{m}$ de turba. (2) $\qty{33000}{mm}$ a
\qty{1}{mm} al año: unos \num{33000} años. (3) La turba solo se conserva bajo
agua estancada, fuera del alcance del oxígeno; a medida que se acumula, el
suelo tiene que hundirse al mismo ritmo para que su superficie siga
encharcada.
\end{solution}

\begin{solution}{exo:g11:fossil-fuel-geology:8}
(1) $20 + 30 \times 1.2 = \qty{56}{\degreeCelsius}$;
$20 + 30 \times 2.3 = \qty{89}{\degreeCelsius}$;
$20 + 30 \times 4.0 = \qty{140}{\degreeCelsius}$.
(2) Hace 120 millones de años, inmadura; hace 95 millones de años, en la
\omterm{def:g11:fossil-fuel-geology:oil-window}{ventana del petróleo}; hace 20 millones de años, en la \omterm{def:g11:fossil-fuel-geology:oil-window}{ventana del gas}.
\end{solution}

\begin{solution}{exo:g11:fossil-fuel-geology:9}
(1) $2580 - 2400 = \qty{180}{m}$. (2) El gas ocupa
$2450 - 2400 = \qty{50}{m}$, y el petróleo, $2580 - 2450 = \qty{130}{m}$.
(3) Se derraman por el punto de derrame, bajo la curva de nivel cerrada más
profunda, y siguen migrando hacia arriba, hacia otra \omterm{def:g11:fossil-fuel-geology:trap}{trampa} o hacia la
superficie.
\end{solution}

\begin{solution}{exo:g11:fossil-fuel-geology:10}
(1) Techo de la ventana a \qty{3}{km}: $(60 - 20)/3 \approx
\qty{13}{\degreeCelsius/km}$; base: $(120 - 20)/3 \approx
\qty{33}{\degreeCelsius/km}$. (2) A \qty{3}{km}, la temperatura va de
$20 + 75 = \qty{95}{\degreeCelsius}$ a $20 + 102 = \qty{122}{\degreeCelsius}$:
en la \omterm{def:g11:fossil-fuel-geology:oil-window}{ventana del petróleo}, o en su base para los gradientes más
pronunciados. La regla es razonable: en la mayoría de las cuencas, una \omterm{def:g11:fossil-fuel-geology:source-rock}{roca
generadora} bajo \qty{3}{km} de rocas ha empezado a producir petróleo.
\end{solution}

\begin{solution}{exo:g11:fossil-fuel-geology:11}
(1) \qty{30}{\percent}: lignito; \qty{70}{\percent}: \omterm{def:g11:fossil-fuel-geology:coal}{carbón} bituminoso;
\qty{92}{\percent}: antracita. (2) La antracita, más allá de unos
\qty{5}{km} y \qty{150}{\degreeCelsius}. (3) Negra y vítrea.
\end{solution}

\begin{solution}{exo:g11:fossil-fuel-geology:12}
(1) $1100/25 = \qty{44}{m}$ por millón de años. (2) $0.044 \times 30
\approx \qty{1.3}{\degreeCelsius}$ por millón de años. (3) La ventana abarca
\qty{60}{\degreeCelsius}: $60/1.32 \approx 45$ millones de años. (4) El
enterramiento se frenó después de hace 95 millones de años (unos 24, luego 23
y \qty{20}{m} por millón de años): la \omterm{def:g11:fossil-fuel-geology:source-rock}{roca generadora} se calentó más despacio.
\end{solution}

\begin{solution}{exo:g11:fossil-fuel-geology:13}
(1) $(80 - 51)/(117 - 51) = 29/66 \approx 0.44$: alrededor del \qty{44}{\percent}.
(2) Migró hacia arriba antes de que existiera encima ninguna \omterm{def:g11:fossil-fuel-geology:trap}{trampa}: escapó a
la superficie, donde se oxidó, o se acumuló en alguna otra \omterm{def:g11:fossil-fuel-geology:trap}{trampa}. (3) Hace 40
millones de años la \omterm{def:g11:fossil-fuel-geology:source-rock}{roca generadora} ya había salido de la \omterm{def:g11:fossil-fuel-geology:oil-window}{ventana del
petróleo}: solo producía gas, y el anticlinal solo habría podido contener gas
procedente de ella.
\end{solution}

\begin{solution}{exo:g11:fossil-fuel-geology:14}
(1) Techo: $(60 - 20)/40 = \qty{1.0}{km}$; base: $(120 - 20)/40 =
\qty{2.5}{km}$. (2) \qty{1.0}{km} está entre \qty{0.4}{km} (hace 150 millones
de años) y \qty{1.2}{km} (hace 120): $150 - 30 \times 0.6/0.8 = 127.5$
millones de años. \qty{2.5}{km} está entre \qty{2.3}{km} (hace 95) y
\qty{3.0}{km} (hace 66): $95 - 29 \times 0.2/0.7 \approx 86.7$ millones de
años. (3) Petróleo no: la \omterm{def:g11:fossil-fuel-geology:source-rock}{roca generadora} salió de la \omterm{def:g11:fossil-fuel-geology:oil-window}{ventana del petróleo}
antes de que se formara el anticlinal, y su petróleo ya había migrado. Podría
contener el gas producido después de hace 80 millones de años.
\end{solution}

\begin{solution}{exo:g11:fossil-fuel-geology:15}
(1) Al menos \qty{5}{km}. (2) $15 + 30 \times 5 = \qty{165}{\degreeCelsius}$,
por encima de \qty{150}{\degreeCelsius}: coherente. (3) Se han levantado y
erosionado al menos \qty{5}{km} de roca. $\qty{5e6}{mm}$ en como mucho
\num{299} millones de años dan al menos $5 \times 10^6/2.99 \times 10^8
\approx \qty{0.017}{mm}$ al año de media.
\end{solution}

\begin{solution}{pb:g11:fossil-fuel-geology:1}
\textbf{1.} \omterm{def:g11:fossil-fuel-geology:source-rock}{Roca generadora}: la lutita negra; almacén: la arenisca porosa;
\omterm{def:g11:fossil-fuel-geology:reservoir}{roca sello}: la fangolita que tiene encima.

\textbf{2.} El plancton que se depositó en un fondo marino pobre en oxígeno
escapó a la descomposición y quedó enterrado en el fango.

\textbf{3.} La \omterm{def:g11:fossil-fuel-geology:source-rock}{roca generadora}, en el Jurásico (165 está entre 201.4 y
143.1). La arenisca empezó en el Jurásico y se depositó sobre todo en el
Cretácico inferior.

\textbf{4.} La materia orgánica dispersa en una roca sedimentaria, insoluble en
los disolventes orgánicos. Solo produce petróleo cuando se calienta hasta unos
\qty{60}{\degreeCelsius}, para lo que hace falta el enterramiento.

\textbf{5.} $400 + 800 + 1100 + 700 + 700 + 300 - 500 = \qty{3500}{m}$.

\textbf{6.} \qty{0.4}{km}, \qty{1.2}{km}, \qty{2.3}{km}, \qty{3.0}{km},
\qty{3.7}{km}, \qty{4.0}{km}; hoy, \qty{3.5}{km}.

\textbf{7.} $400/15 \approx 26.7$; $800/30 \approx 26.7$; $1100/25 = 44$;
$700/29 \approx 24.1$; $700/31 \approx 22.6$; $300/15 = 20$ (\unit{m} por
millón de años).

\textbf{8.} Hace 20 millones de años, a \qty{4.0}{km}.

\textbf{9.} La cuenca se levantó, y la erosión retiró \qty{500}{m} de roca de
su parte superior.

\textbf{10.} La base de la arenisca se depositó hace 150 millones de años,
sobre \qty{400}{m} de lutita y fangolita: $3.5 - 0.4 = \qty{3.1}{km}$.

\textbf{11.} $T = 20 + 30 z$: \qty{32}{\degreeCelsius}, \num{56}, \num{89},
\num{110}, \num{131}, \qty{140}{\degreeCelsius}; hoy,
\qty{125}{\degreeCelsius}.

\textbf{12.} $40/30 \approx \qty{1.33}{km}$ y $100/30 \approx
\qty{3.33}{km}$.

\textbf{13.} Entre hace 120 millones de años (\qty{56}{\degreeCelsius}) y
hace 95 millones de años (\qty{89}{\degreeCelsius}).

\textbf{14.} La temperatura subió 33 grados en 25 millones de años; tardó
$25 \times 4/33 \approx 3.0$ millones de años en subir los 4 grados que van de
\num{56} a \qty{60}{\degreeCelsius}: hace unos $120 - 3.0 = 117$ millones de
años.

\textbf{15.} Entre hace 66 millones de años (\qty{110}{\degreeCelsius}) y
hace 35 (\qty{131}{\degreeCelsius}): hace $66 - 31 \times 10/21 \approx 51.2$
millones de años.

\textbf{16.} \qty{140}{\degreeCelsius}, hace 20 millones de años. A
\qty{125}{\degreeCelsius}, hoy está en la \omterm{def:g11:fossil-fuel-geology:oil-window}{ventana del gas}: produce sobre todo
metano.

\textbf{17.} Sin ninguna \omterm{def:g11:fossil-fuel-geology:trap}{trampa} encima, migró hacia arriba y escapó, a la
superficie o a otras \omterm{def:g11:fossil-fuel-geology:trap}{trampas}.

\textbf{18.} $(80 - 51.2)/(117.0 - 51.2) \approx 0.44$.

\textbf{19.} Gas arriba, luego petróleo, luego agua.

\textbf{20.} La \omterm{def:g11:fossil-fuel-geology:source-rock}{roca generadora} entró en la \omterm{def:g11:fossil-fuel-geology:oil-window}{ventana del petróleo} hace unos
\textbf{117 millones de años}.
\end{solution}
